The midterm results already answer RQ1 for GEARS (it does not beat addition) and reduce RQ2
for GEARS to a comparison of two additive constructions.

\noindent\textbf{$\eta$ is reproducible across cells.} The worry was that ``highly
non-additive'' pairs might just be the noisiest ones, making their high error a
regression-to-the-mean artefact. The split-half design rules this out: $\eta$ replicates
across disjoint cells (Spearman 0.94), and every error is scored on the half not used to
define the groups. Reproducibility alone does not show that $\eta$ measures interaction rather than perturbation
size; the caveat at the end of this section addresses that. The high-$\eta$ tertile is where a
model with real interaction knowledge has room to beat addition.

\noindent\textbf{Addition is a hard bar, and the room above it is concentrated.} Most of
addition's error sits in the top $\eta$ tertile, five times the low tertile's. An overall
average therefore mostly reflects easy, near-additive pairs, so the rollout analysis will
report advantage over addition as a function of $\eta$, not as one average.

\noindent\textbf{Ablation: smoothing the singles does not help.} Low-rank addition is worse
than plain addition on 82\% of pairs. With pseudobulks from roughly 100 to 300 cells, the
single effects are already precise, so the rank-10 projection removes more real signal than
noise. A singles-only model must know something about how effects interact to beat addition.

\noindent\textbf{GEARS has not learned interactions, and it cannot carry state.} On the
strict split, GEARS's error exceeds addition's in every tertile; its gap narrows in the high
tertile (0.194 vs.\ 0.153) but does not reverse. Its predicted shift also ignores the input
cell, so its rollouts are fixed by construction:
\begin{equation}
\text{carried: } x + g(A) + g(B), \qquad \text{restated: } x + g(A) + g(A,B).
\end{equation}
The carried rollout cannot depend on order, since $g(A) + g(B)$ is symmetric. The
restated rollout can: it equals the joint prediction plus an extra $g(A)$, so it counts $A$
twice and should overshoot GEARS's joint prediction by exactly $g(A)$. Applying $A$ twice
simply doubles its shift. For GEARS, the remaining empirical question is whether its own
singles, $g(A) + g(B)$, land closer to the measured double than its joint $g(A,B)$. Because
$g(A)$ is trained to match each measured single, we expect the carried rollout to score
close to the \addnull, and so better than GEARS's joint prediction; that would reflect
addition, not composition. scGPT takes input expression by design, so composition is an open question for it; the
round-trip gate will test whether its output actually depends on its input.

\noindent\textbf{Caveat: $\eta$ partly tracks effect size.} $\eta$ correlates with the size
of the double's shift (Spearman 0.54), and shift size alone predicts additive error about as
well (0.58). $\eta$ relative to the shift barely predicts it (0.15), and the scale-free \pd
hardly falls across tertiles (0.880 to 0.862). Part of the trend is ``bigger perturbations
have bigger absolute errors''. For the final report we will compare models with addition
within tertiles, which cancels scale, and check conclusions with relative $\eta$ or effect
size as a covariate.
